Consider two explanations for the same two-time record. In the first, the source is asymmetric and the sensors add no delay. In the second, the source is reversible and one sensor delays it. Each channel separately has variance one and no covariance between distinct times, called a unit white marginal spectrum. The first explanation has cross spectrum
\[
 F^A_{12}=\tfrac1{10}+\tfrac15e^{-i\omega}.
\]
A reversible source with $S^R_{12}=\tfrac15+\tfrac15\cos\omega$, observed after delaying channel 1 by one sample, instead gives
\[
 F^R_{12}=e^{-i\omega}S^R_{12}
 =\tfrac1{10}+\tfrac15e^{-i\omega}+\tfrac1{10}e^{-2i\omega}.
\]
The spectral eigenvalues are at least $7/10$ and $3/5$, respectively, so both constructions are positive and define valid Gaussian processes. Their cross-covariances are
\[
\begin{array}{c|ccccc}
 k&-2&-1&0&1&2\\\hline
 C^A_{12}(k)&0&0&1/10&1/5&0\\
 C^R_{12}(k)&0&0&1/10&1/5&1/10
\end{array}
\]
A two-time record sees only lags $-1,0,1$, the middle three columns. Its complete covariance, and hence its Gaussian law, agrees under both explanations. The extra term in the second spectrum changes only the last column shown. Lag two would distinguish the explanations, but is absent from that record.
